\documentclass{jsarticle}
\usepackage[dvipdfmx,final]{graphicx}
\usepackage{amsmath}
\usepackage{amssymb}
\begin{document}
\title{Gamma function relativity restream into Mebius \\
Klein dimension, Cohomological define with Euler product, \\
Euler number from integral manifold to general relativity \\
and Special relativity. Imaginary pole of Volume manifold \\
emerged with Beta function.}
\author{Masaaki Yamaguchi}
\date{}
\maketitle

        オイラー数を多様体積分で施すと、
   ガンマ関数になる。この場合は、Frobeniusの定理によって、
   和と積の法則、確率論から導かれる結果、
   異次元と宇宙の関係に持っていける。
   オイラー数を多様体積分すると、ガンマ関数になる。
   オイラー数における頂点が、宇宙の特異点に対応して、
   辺がノルムに属して、面が一般相対性理論における無限の接線であり、
   $$ \text{F(頂点) - N(辺) + E(面)} = 2 $$
   $$ {(x,y,z) \over \Gamma} = \int{(- F\circ N + E)} + 
   (-F\circ E + N) + (N\circ E + F)dx_m $$
   $$ {{{\partial \over \partial f_S}(x,y,z)} \over {\Gamma}} $$
   ここで、オイラーの定数をも多様体積分にして、付け加えると、
   $$ = \int{({\int {1 \over x^s}dx} - \log x)}\mathrm{dvol} $$
   大域的多様体積分のガンマ関数に化ける。
   $$ = \int e^{-x}x^{1 - t}dx_m $$

   $AdS_5$多様体の方程式を、オイラー数の各部分単体に合わせると、
   $$ ||ds^2|| = e^{-2\pi T||\psi||}[\eta + \bar{h}(x)]dx^{\mu}dx^{\nu} 
   + T^2d^2 \psi $$
   開集合として、表せる。
   $$ = \mathcal{O}(x) $$
   この上の式たちから、ダランベルシアンは、$\square$と、オイラー数から演算子
   が決まっていて、共変微分は、$\nabla$で、双対被覆での単体での演算子で、
   決まっている。
   $$ \square = {{8 \pi G} \over c^4} T^{\mu\nu} $$
   $$ \nabla \psi = 4 \pi G \rho $$

   オイラー数を共変微分すると、偏微分方程式と同じ機能を、
   大域的積分多様体のオイラー数が、受け持っていて、
   このオイラー数の大域的多様体の共変微分で、ガンマ関数の大域的多様体
   が、
   ここで、オイラーの定数をも多様体積分にして、付け加えると、
   $$ = \int{({\int {1 \over x^s}dx} - \log x)}\mathrm{dvol} $$
   大域的多様体積分のガンマ関数に化ける。
   $$ = \int e^{-x}x^{1 - t}dx_m $$
   $$ = {d \over d f}F = \int \Gamma(\gamma)^{'}dx_m $$
   $$ = \int C dx_m = \int \Gamma dx_m \circ {d \over d \gamma}\Gamma $$
   $$ \le \int \Gamma dx_m + {d \over d \gamma}\Gamma $$
   $$ = e^{- \theta} + e^{i \theta} $$
   $$ = \square = 2(\cos (i x \log x) - i \sin (i x \log x))$$
   $$ = e^{-f} - e^{f} \le e^{-f} + e^{f} $$
   統一場理論は、彩さんのJones多項式に行き着く。

   Zata function escourt into Forgotfull of underlying modify on Sum
   of summative group, thiis Forgotfull summative group instented with
   Space ideal theorem from quantum level of deprivate space in aspect
   of quantum level of differential geometry,
   This theorem construct with Higgs field and Morse theory, moreover
   Hortshorn conjecture, fourth step deduction of theorem from 
   Euler product of integral manifold.
   This manfold revolte with global differential manifold in global topology.
   From these theorem inspirate from Forgotfull theorem in Space ideality
   theory, and this theory conclude into entropy of non exchange equation.
   Forgotten theory estimate with zeta function oneselves.
   Forgotten theory income with non relativity theory and this theory
   face to face without partner and partner of non relativity from
   ideal of space in revolution theory.
   All exist and non exist partner with non relativity escourt into 
   Forgotten theory, and this theory include with same underly group
   of all of partner with same distance ideal.

   $$ X \dotsi \to Y, Y \dotsi \to X $$
   $$ \bigoplus(i \hbar^{\nabla})^{\oplus L} = e^{-x \log x} $$
   This equation of quantum level of differential geomerty
   instimate with entorpy of non exchange equation.
   This equation means with low level of botton entropy in Space ideal of
   Forgotten theory. Moveover, this equation enterstein with all of exist
   theory from general relativity.
   No time and No Space of relativity, and monotonicity of magnetic component
   with gravity and antigravity, and this theory comment with partner of
   magnetic component theory included with being resulted from Forgotten
   theory.
   $$ \square \to \nabla |: x \to y, f(x) \to g(y), f^{-1}(x)xf(x) 
   \dotsi \to g(y) $$
   $$ \square \dotsi \to \nabla |: f(x) \cong g(y) $$
   $$ {d \over dt}g_{ij}(t) = - 2 R_{ij}, F^m_t = -2\int{{(R + \nabla_i \nabla
   _j f)} \over {(\Delta + f)}}dm $$
   $$ {d \over df}F = m(x) $$
   $$ F|_{\_} = {d \over df}F, F^m_t = e^{-f}dV $$
   $$ \bigoplus(i \hbar^{\nabla})^{\oplus L} = e^{-x \log x} $$
   These equation represent with Forgotten theory.
   空間概念の量子化は、代数幾何の量子化であり、忘却関数でもあり、
   忘却関手のカテゴリー論でもあり、簡潔に言うと、すべての相対性が
   機能すると、全て同型というコホモロジー論でのコイコライザーである。
   このコイコライザーが、空間概念の量子化である。
   すべては、ゼータ関数へと行き着く。
   ベータ関数は、誤差関数であり、宇宙の雑音として存在していると、
   異次元へと移行する。
   $$ \beta(p,q) \cong {\Gamma(p)\Gamma(q) \le \Gamma(p+q)}, 
   {{\Gamma(p)\Gamma(q)} \over \Gamma(p+q)} \le 1 $$
   $$ \int e^{-x}x^{1-t}dx_m = \int e^{-x}x^{1-t}dx\cdot e^{-\theta} $$
   $$ = \Gamma^{{\gamma}^{'}} $$
   $$ = e^{-x \log x} $$
   閉３次元多様体と３次元球面が、空間概念の量子化としての結果であり、
   この多様体が、無としてのベータ関数を形作っている３次元多様体である。
   宇宙と異次元合わせてが、空間概念の量子化での単体である。
   コイコライザーとは、共変微分によって対になっている変数に作用する写像であり、
   このコイコライザーで処理された変数同士を同型にさせる群が、
   同型のコイコライザーである。そして、最小の単体でもある。
   故に、宇宙と異次元合わせて、コイコライザーである。
   球対称である一般相対性理論の平坦な空間が、空間概念の量子化でもある。
   ビッグ・バンは、イコライザーであり、
   $$ \langle f,g\rangle \mapsto \langle d,d\rangle $$
   として、宇宙と異次元の曲平面を表している。
 
   [Reference: 深谷賢治先生と加藤文元先生と竹内薫先生、S.マックレーン先生]



\end{document}
